\documentclass[11pt, a4paper]{article}

\usepackage{amsmath, amssymb, amsthm}
\usepackage{geometry}
\usepackage{booktabs}
\usepackage{array}
\usepackage{parskip}
\usepackage{microtype}
\usepackage{hyperref}
\usepackage{enumitem}
\usepackage{xcolor}
\usepackage{mdframed}

\geometry{margin=2.5cm, top=2.8cm, bottom=2.8cm}

\hypersetup{
    colorlinks=true,
    linkcolor=black,
    urlcolor=black,
}

% Exercise environment
\newcounter{exercise}[section]
\renewcommand{\theexercise}{\thesection.\arabic{exercise}}

\newenvironment{exercise}[1][]{%
    \refstepcounter{exercise}%
    \par\medskip\noindent%
    \textbf{Exercise~\theexercise}%
    \ifx\relax#1\relax\else\ ---\ #1\fi\par\smallskip%
}{%
    \par\medskip%
}

% Motivation box
\newmdenv[
    linewidth=0.6pt,
    linecolor=black!40,
    backgroundcolor=black!3,
    innerleftmargin=10pt,
    innerrightmargin=10pt,
    innertopmargin=6pt,
    innerbottommargin=6pt,
    skipabove=\medskipamount,
    skipbelow=\medskipamount,
]{motivbox}

% Solution environment (for solutions section)
\newcounter{solution}[section]
\renewcommand{\thesolution}{\thesection.\arabic{solution}}

\newenvironment{solution}[1][]{%
    \refstepcounter{solution}%
    \par\medskip\noindent%
    \textbf{Solution~\thesolution}%
    \ifx\relax#1\relax\else\ ---\ #1\fi\par\smallskip%
}{%
    \par\medskip%
}

% Inline code
\newcommand{\code}[1]{\texttt{#1}}

% Convenience: bold vectors
\renewcommand{\vec}[1]{\mathbf{#1}}

\title{\textbf{Linear Algebra for Mechanistic Interpretability}\\[0.4em]
\large Pencil-and-Paper Exercises}
\author{}
\date{}

\begin{document}

\maketitle
\thispagestyle{empty}

\begin{abstract}
\noindent
\textbf{Who this is for.} Someone with a CS background whose linear algebra is rusty,
working toward understanding how transformer attention works mechanically.

\textbf{How to use this.} Work through each section in order. The exercises build on each
other. Do them by hand --- the point is to build intuition, not to get answers fast.
Check your work against the solutions starting on page~\pageref{sec:solutions}.
\end{abstract}

\tableofcontents
\newpage

% ──────────────────────────────────────────────────────────────────────────────
\section{Vectors and Dot Products}
% ──────────────────────────────────────────────────────────────────────────────

\begin{motivbox}
\textbf{Why this matters for transformers.}
Every token in a transformer is represented as a vector (its ``embedding''). When the model
computes attention, it is fundamentally asking: \emph{how similar is this query vector to
each key vector?} That similarity is measured with a dot product. If you can read a dot
product and know what it means geometrically, you can read an attention pattern and know
what the model is ``looking at.''
\end{motivbox}

\begin{exercise}[Basic dot product (2D)]
Compute the dot product of:
\[
    \vec{a} = \begin{bmatrix} 3 \\ 1 \end{bmatrix}, \qquad
    \vec{b} = \begin{bmatrix} 2 \\ 4 \end{bmatrix}
\]
Recall: $\vec{a} \cdot \vec{b} = a_1 b_1 + a_2 b_2$.
\end{exercise}

\begin{exercise}[Basic dot product (3D)]
Compute:
\[
    \vec{u} \cdot \vec{v} = \begin{bmatrix} 1 \\ -2 \\ 3 \end{bmatrix} \cdot
                            \begin{bmatrix} 4 \\ 1 \\ 2 \end{bmatrix}
\]
\end{exercise}

\begin{exercise}[Geometric interpretation]
You are told that $\vec{a} \cdot \vec{b} = \|\vec{a}\| \|\vec{b}\| \cos\theta$,
where $\theta$ is the angle between the vectors.

For each case below, say whether the dot product will be \textbf{positive},
\textbf{negative}, or \textbf{zero} --- and explain in one sentence what that
means geometrically:
\begin{enumerate}[label=(\alph*)]
    \item Two vectors pointing in exactly the same direction.
    \item Two vectors pointing in exactly opposite directions.
    \item Two vectors that are perpendicular (orthogonal).
    \item A long vector and a short vector pointing in the same direction, compared to
          two unit vectors pointing in the same direction.
          Which dot product is larger?
\end{enumerate}
\end{exercise}

\begin{exercise}[Dot product as similarity]
Consider three 2-dimensional word embedding vectors (toy, made-up numbers):
\[
    \text{``cat''} = \begin{bmatrix} 2 \\ 3 \end{bmatrix}, \qquad
    \text{``dog''} = \begin{bmatrix} 2 \\ 2 \end{bmatrix}, \qquad
    \text{``car''} = \begin{bmatrix} -3 \\ 1 \end{bmatrix}
\]
\begin{enumerate}[label=(\alph*)]
    \item Compute $\text{``cat''} \cdot \text{``dog''}$.
    \item Compute $\text{``cat''} \cdot \text{``car''}$.
    \item Which pair is more ``similar'' according to the dot product?
          Does this match your intuition about the words?
\end{enumerate}
\end{exercise}

\begin{exercise}[Normalisation and dot products]
Two vectors:
\[
    \vec{p} = \begin{bmatrix} 3 \\ 4 \end{bmatrix}, \qquad
    \vec{q} = \begin{bmatrix} 0 \\ 5 \end{bmatrix}
\]
\begin{enumerate}[label=(\alph*)]
    \item Compute $\|\vec{p}\| = \sqrt{p_1^2 + p_2^2}$ and $\|\vec{q}\|$.
    \item Compute $\vec{p} \cdot \vec{q}$.
    \item Normalise both vectors (divide each by its magnitude) to get
          $\hat{\vec{p}}$ and $\hat{\vec{q}}$.
    \item Compute $\hat{\vec{p}} \cdot \hat{\vec{q}}$. This is the \emph{cosine
          similarity}. Why might cosine similarity be a more useful measure of
          directional similarity than the raw dot product?
\end{enumerate}
\end{exercise}

% ──────────────────────────────────────────────────────────────────────────────
\section{Matrix Multiplication}
% ──────────────────────────────────────────────────────────────────────────────

\begin{motivbox}
\textbf{Why this matters for transformers.}
The weight matrices $W_Q$, $W_K$, $W_V$, and $W_O$ are all matrices. When a token
embedding passes through one of these, it is a matrix multiplication. Understanding what
matrix multiplication does --- and critically, what \emph{shape} the result has --- is
essential for reading transformer code without getting lost.
\end{motivbox}

\begin{exercise}[$2 \times 2$ matrix multiplication]
Compute $AB$ where:
\[
    A = \begin{bmatrix} 1 & 2 \\ 3 & 4 \end{bmatrix}, \qquad
    B = \begin{bmatrix} 5 & 6 \\ 7 & 8 \end{bmatrix}
\]
Recall: the entry in row $i$, column $j$ of $AB$ is the dot product of row $i$ of $A$
with column $j$ of $B$.
\end{exercise}

\begin{exercise}[Non-square matrix multiplication]
Compute $AB$ where:
\[
    A = \begin{bmatrix} 1 & 0 & 2 \\ 3 & 1 & 1 \end{bmatrix} \quad (2 \times 3), \qquad
    B = \begin{bmatrix} 1 & 2 \\ 0 & 1 \\ 4 & 0 \end{bmatrix} \quad (3 \times 2)
\]
\begin{enumerate}[label=(\alph*)]
    \item What shape is $AB$?
    \item Compute $AB$.
    \item Can you compute $BA$? If yes, what shape is it? If no, why not?
\end{enumerate}
\end{exercise}

\begin{exercise}[Shape tracking]
For each of the following, state whether the multiplication is valid and, if so,
give the shape of the result. You do \textbf{not} need to compute the values.
\begin{enumerate}[label=(\alph*)]
    \item $A$ is $(3 \times 4)$, $B$ is $(4 \times 2)$. Compute $AB$.
    \item $A$ is $(5 \times 3)$, $B$ is $(5 \times 3)$. Compute $AB$.
    \item $A$ is $(1 \times 6)$ (a row vector), $B$ is $(6 \times 1)$ (a column vector).
          Compute $AB$.
    \item $A$ is $(6 \times 1)$, $B$ is $(1 \times 6)$. Compute $AB$.
          How is this different from~(c)?
    \item In a transformer, a token embedding $\vec{x}$ has shape $(d_{\mathrm{model}},)$,
          treated as $(1 \times d_{\mathrm{model}})$. The query weight matrix $W_Q$ has
          shape $(d_{\mathrm{model}} \times d_{\mathrm{head}})$.
          What shape is $\vec{x} W_Q$?
\end{enumerate}
\end{exercise}

\begin{exercise}[$W_V$ and $W_O$ in miniature]
In a transformer attention head, $W_V$ maps each token to a ``value'' vector, and
$W_O$ maps the aggregated value back to the residual stream. Here is a tiny version.

Let:
\[
    W_V = \begin{bmatrix} 1 & 0 \\ 0 & 2 \end{bmatrix}, \qquad
    W_O = \begin{bmatrix} 0 & 1 \\ 1 & 0 \end{bmatrix}
\]
and let token embedding $\vec{x} = \begin{bmatrix} 3 & 1 \end{bmatrix}$ (a row vector).
\begin{enumerate}[label=(\alph*)]
    \item Compute $\vec{v} = \vec{x} W_V$ (the value vector for this token).
    \item Compute the output $\vec{o} = \vec{v} W_O$.
    \item What is the combined effect $\vec{x}(W_V W_O)$?
          Compute $W_V W_O$ first, then multiply by $\vec{x}$.
    \item Does your answer to~(c) match~(b)? It should --- this is matrix
          associativity: $(AB)C = A(BC)$.
\end{enumerate}
\end{exercise}

\begin{exercise}[The OV circuit]
This exercise previews one of the key ideas in mechanistic interpretability:
the \textbf{OV circuit}, which describes what a head \emph{does} to the residual
stream regardless of what it attends to.

The OV matrix is defined as $W_{OV} = W_V W_O$. Using:
\[
    W_V = \begin{bmatrix} 2 & 1 \\ 0 & 3 \end{bmatrix}, \qquad
    W_O = \begin{bmatrix} 1 & 0 \\ 0 & 1 \end{bmatrix}
\]
\begin{enumerate}[label=(\alph*)]
    \item Compute $W_{OV} = W_V W_O$.
    \item What does $W_O = I$ (the identity matrix) imply about what $W_O$ is doing?
    \item Apply $W_{OV}$ to two different token embeddings:
          $\vec{x}_1 = \begin{bmatrix} 1 & 0 \end{bmatrix}$ and
          $\vec{x}_2 = \begin{bmatrix} 0 & 1 \end{bmatrix}$.
          What do you notice about the relationship between input and output?
\end{enumerate}
\end{exercise}

% ──────────────────────────────────────────────────────────────────────────────
\section{Linear Transformations}
% ──────────────────────────────────────────────────────────────────────────────

\begin{motivbox}
\textbf{Why this matters for transformers.}
A matrix is not just an array of numbers --- it is a \emph{function}. When you multiply
a vector by a matrix, you are applying a linear transformation: stretching, rotating, or
projecting that vector in space. Every weight matrix in a transformer is doing this to
token representations. If you can think of matrices as functions, you can think about
what information they extract or suppress.
\end{motivbox}

\begin{exercise}[Matrix as function]
Let $T = \begin{bmatrix} 2 & 0 \\ 0 & 3 \end{bmatrix}$.

Apply $T$ to each of the following vectors and describe what geometric effect $T$ has:
\begin{enumerate}[label=(\alph*)]
    \item $\vec{v} = \begin{bmatrix} 1 \\ 0 \end{bmatrix}$
    \item $\vec{v} = \begin{bmatrix} 0 \\ 1 \end{bmatrix}$
    \item $\vec{v} = \begin{bmatrix} 1 \\ 1 \end{bmatrix}$
    \item In one sentence, what does $T$ do to any 2D vector?
\end{enumerate}
\end{exercise}

\begin{exercise}[Two matrices, same input]
Let:
\[
    A = \begin{bmatrix} 1 & 0 \\ 0 & -1 \end{bmatrix}, \qquad
    B = \begin{bmatrix} 0 & 1 \\ 1 & 0 \end{bmatrix}, \qquad
    \vec{v} = \begin{bmatrix} 3 \\ 2 \end{bmatrix}
\]
\begin{enumerate}[label=(\alph*)]
    \item Compute $A\vec{v}$. What does $A$ do geometrically?
    \item Compute $B\vec{v}$. What does $B$ do geometrically?
    \item These are two different ``lenses'' through which to view the same vector
          $\vec{v}$. In a transformer, different attention heads apply different
          $W_Q, W_K, W_V$ matrices to the same token --- they are asking different
          questions about the same information. Does this analogy feel concrete now?
\end{enumerate}
\end{exercise}

\begin{exercise}[Matrix multiplication as function composition]
Let $f(\vec{v}) = A\vec{v}$ and $g(\vec{v}) = B\vec{v}$ where:
\[
    A = \begin{bmatrix} 1 & 1 \\ 0 & 1 \end{bmatrix}, \qquad
    B = \begin{bmatrix} 2 & 0 \\ 0 & 2 \end{bmatrix}, \qquad
    \vec{v} = \begin{bmatrix} 1 \\ 2 \end{bmatrix}
\]
\begin{enumerate}[label=(\alph*)]
    \item Compute $f(\vec{v}) = A\vec{v}$.
    \item Compute $g(f(\vec{v})) = B(A\vec{v})$ by applying $B$ to your answer from~(a).
    \item Now compute $C = BA$ directly.
    \item Compute $C\vec{v}$.
    \item Do (b) and (d) agree? They should --- matrix multiplication \emph{is} function
          composition: applying $A$ then $B$ is identical to applying the single matrix
          $BA$.
\end{enumerate}
\end{exercise}

\begin{exercise}[Projections]
A projection matrix onto the $x$-axis is:
\[
    P = \begin{bmatrix} 1 & 0 \\ 0 & 0 \end{bmatrix}
\]
\begin{enumerate}[label=(\alph*)]
    \item Apply $P$ to $\vec{v} = \begin{bmatrix} 4 \\ 7 \end{bmatrix}$.
          What information is lost?
    \item Apply $P$ twice: compute $P(P\vec{v})$. What do you notice?
    \item A matrix $M$ satisfying $M^2 = M$ is called \textbf{idempotent}. Verify that
          $P^2 = P$. Why does this make sense geometrically?
          (Projecting something already projected should do nothing.)
    \item In transformers, the residual stream means that information is \emph{added to}
          (not replaced by) the stream. How is this different from applying a projection?
\end{enumerate}
\end{exercise}

% ──────────────────────────────────────────────────────────────────────────────
\section{Attention Scores by Hand}
% ──────────────────────────────────────────────────────────────────────────────

\begin{motivbox}
\textbf{Why this matters.}
This is the core computation. In a transformer attention head, each token generates a
\textbf{query} (what am I looking for?) and a \textbf{key} (what do I offer?). The
attention score between token $i$ and token $j$ is the dot product of token $i$'s query
with token $j$'s key, scaled and softmaxed. Working through this by hand makes the
formula completely concrete.

The full formula for attention is:
\[
    \mathrm{Attention}(Q, K, V) =
    \mathrm{softmax}\!\left(\frac{QK^T}{\sqrt{d_{\mathrm{head}}}}\right) V
\]
In this section you will compute the $QK^T / \sqrt{d_{\mathrm{head}}}$ part and the
softmax. We use $d_{\mathrm{head}} = 2$ throughout.
\end{motivbox}

\begin{exercise}[Computing Q and K vectors]
Suppose we have a 3-token sequence. The token embeddings (already 2D for simplicity) are:
\[
    \vec{x}_1 = \begin{bmatrix} 1 \\ 0 \end{bmatrix}, \qquad
    \vec{x}_2 = \begin{bmatrix} 0 \\ 1 \end{bmatrix}, \qquad
    \vec{x}_3 = \begin{bmatrix} 1 \\ 1 \end{bmatrix}
\]
The query and key weight matrices are:
\[
    W_Q = \begin{bmatrix} 1 & 0 \\ 0 & 1 \end{bmatrix}, \qquad
    W_K = \begin{bmatrix} 2 & 0 \\ 0 & 1 \end{bmatrix}
\]
\begin{enumerate}[label=(\alph*)]
    \item Compute the query vectors $\vec{q}_i = \vec{x}_i^T W_Q$ for $i = 1, 2, 3$.
    \item Compute the key vectors $\vec{k}_i = \vec{x}_i^T W_K$ for $i = 1, 2, 3$.
\end{enumerate}
\textit{Note:} We write embeddings as row vectors and multiply $\vec{x}^T W$ rather than
$W^T \vec{x}$ --- this matches the convention used in ARENA and most transformer
implementations.
\end{exercise}

\begin{exercise}[Raw attention scores]
Using your $Q$ and $K$ vectors from Exercise~4.1:
\begin{enumerate}[label=(\alph*)]
    \item Compute the raw attention score $s_{ij} = \vec{q}_i \cdot \vec{k}_j$ for every
          pair $(i, j)$ where $i, j \in \{1, 2, 3\}$. Fill in this $3 \times 3$ table:

\medskip
\begin{center}
\begin{tabular}{c|ccc}
\toprule
          & $\vec{k}_1$ & $\vec{k}_2$ & $\vec{k}_3$ \\
\midrule
$\vec{q}_1$ & $s_{11}$ & $s_{12}$ & $s_{13}$ \\
$\vec{q}_2$ & $s_{21}$ & $s_{22}$ & $s_{23}$ \\
$\vec{q}_3$ & $s_{31}$ & $s_{32}$ & $s_{33}$ \\
\bottomrule
\end{tabular}
\end{center}
\medskip

    \item Write the full attention score matrix $S$ ($3 \times 3$).
    \item Row $i$ of $S$ contains the scores for token $i$ attending to every other
          token. Which token does token~3 most ``want to attend to'' based on raw scores?
\end{enumerate}
\end{exercise}

\begin{exercise}[Scaling and softmax]
\textbf{Step 1: Scale.}
Divide every entry in your score matrix $S$ by
$\sqrt{d_{\mathrm{head}}} = \sqrt{2} \approx 1.414$.
Call the result $\tilde{S}$. (Round to 2 decimal places.)

\medskip
\textbf{Step 2: Softmax.}
The softmax of a row vector $\vec{z} = [z_1, z_2, \ldots, z_n]$ is:
\[
    \mathrm{softmax}(\vec{z})_i = \frac{e^{z_i}}{\displaystyle\sum_j e^{z_j}}
\]
Apply softmax to each row of $\tilde{S}$ independently. This gives the
\textbf{attention pattern} $A$ --- a $3 \times 3$ matrix where each row sums to~1.

\medskip
Useful values:
$e^0 \approx 1.000$,\;
$e^{0.71} \approx 2.034$,\;
$e^{1.41} \approx 4.096$,\;
$e^{2.12} \approx 8.331$.

\begin{enumerate}[label=(\alph*)]
    \item Compute the scaled scores for row~1 (token~1's query against all keys).
    \item Apply softmax to row~1. Verify it sums to~1.
    \item Apply softmax to row~3 (you can leave row~2 as an exercise if time is short).
    \item Interpret the attention pattern: what is token~1 ``attending to''?
          What is token~3 ``attending to''?
\end{enumerate}
\end{exercise}

% ──────────────────────────────────────────────────────────────────────────────
\section{Tensor Shapes and Einsum Notation}
% ──────────────────────────────────────────────────────────────────────────────

\begin{motivbox}
\textbf{Why this matters.}
Modern transformer code uses \code{einops.einsum} and \code{torch.einsum} extensively
because they make the intention of each operation explicit via index notation. Being able
to read an einsum string and immediately know the input/output shapes --- and what
computation is happening --- is a superpower when reading or debugging transformer code.
The ARENA curriculum uses these constantly.

\textbf{How to read einsum notation.} Each letter is a dimension. A letter appearing in
both inputs is \emph{summed over} (contracted). A letter appearing in the output is
\emph{kept}. Shapes must be consistent: the same letter must have the same size everywhere
it appears.
\end{motivbox}

\begin{exercise}[Basic einsum interpretation]
For each einsum expression, state: (i) what each letter represents, (ii) the shapes of
the inputs, (iii) the shape of the output, and (iv) in plain English, what operation is
being performed.

\begin{enumerate}[label=(\alph*)]
    \item \code{"i, i -> "} applied to tensors \code{A} and \code{B} where \code{A}
          has shape \code{(4,)} and \code{B} has shape \code{(4,)}.

    \item \code{"i j, j k -> i k"} where \code{A} has shape \code{(3, 4)} and \code{B}
          has shape \code{(4, 5)}.

    \item \code{"b i j, b j k -> b i k"} where \code{A} has shape \code{(8, 3, 4)} and
          \code{B} has shape \code{(8, 4, 5)}.
          (\code{b} is the batch dimension.)

    \item \code{"i j -> j i"} where \code{A} has shape \code{(3, 7)}.
\end{enumerate}
\end{exercise}

\begin{exercise}[Shape consistency check]
For each of the following, identify the error (if any) in the einsum expression given
the tensor shapes:
\begin{enumerate}[label=(\alph*)]
    \item \code{"i j, i k -> j k"} with \code{A} shape \code{(3, 4)} and \code{B}
          shape \code{(3, 5)}.
    \item \code{"i j, j k -> i k"} with \code{A} shape \code{(3, 4)} and \code{B}
          shape \code{(5, 2)}.
    \item \code{"i j k, i j -> i k"} with \code{A} shape \code{(2, 3, 4)} and \code{B}
          shape \code{(2, 3)}.
\end{enumerate}
\end{exercise}

\begin{exercise}[The ARENA 1.3 einsum]
The following einsum appears in ARENA exercise~1.3 for computing attention scores:
\begin{center}
\code{einops.einsum(Q, K, "seqQ n h, seqK n h -> n seqQ seqK")}
\end{center}
where the variables have these meanings:
\begin{itemize}
    \item \code{seqQ}: number of query positions (sequence length)
    \item \code{seqK}: number of key positions (same value as \code{seqQ} in
          self-attention)
    \item \code{n}: number of attention heads
    \item \code{h}: head dimension ($d_{\mathrm{head}}$)
\end{itemize}
\begin{enumerate}[label=(\alph*)]
    \item What are the shapes of \code{Q} and \code{K}?
    \item What is the shape of the output?
    \item Which dimension is being summed over (contracted)? What does summing over
          \code{h} correspond to mathematically?
    \item What does entry \code{[n, i, j]} of the output tensor represent?
    \item Why is \code{n} in the output rather than contracted? What would it mean if we
          had written \code{"seqQ n h, seqK n h -> seqQ seqK"} instead?
\end{enumerate}
\end{exercise}

\begin{exercise}[Writing your own einsum]
Write the einsum string for each of the following operations:
\begin{enumerate}[label=(\alph*)]
    \item Standard (non-batched) matrix multiplication of \code{A} (shape $m \times n$)
          with \code{B} (shape $n \times p$) to get output of shape $m \times p$.

    \item For each token in a sequence, compute the dot product of that token's vector
          with a fixed vector \code{w}. Input: \code{X} of shape \code{(seq, d)},
          \code{w} of shape \code{(d,)}. Output shape: \code{(seq,)}.

    \item Compute the outer product of two vectors \code{u} of shape \code{(m,)} and
          \code{v} of shape \code{(n,)} to get a matrix of shape \code{(m, n)}.

    \item \textbf{Challenge.} You have value vectors \code{V} of shape
          \code{(seqK, n, h)} and an attention pattern \code{A} of shape
          \code{(n, seqQ, seqK)}. Write the einsum that computes the weighted sum of
          values for each query position and head. Output shape: \code{(seqQ, n, h)}.
\end{enumerate}
\end{exercise}

% ──────────────────────────────────────────────────────────────────────────────
\newpage
\section*{Worked Solutions}
\label{sec:solutions}
\addcontentsline{toc}{section}{Worked Solutions}
% ──────────────────────────────────────────────────────────────────────────────

\setcounter{section}{1}
\setcounter{solution}{0}

% --- Section 1 Solutions ---
\subsection*{Section 1: Vectors and Dot Products}

\begin{solution}[Basic dot product (2D)]
\[
    \vec{a} \cdot \vec{b} = (3)(2) + (1)(4) = 6 + 4 = \mathbf{10}
\]
\end{solution}

\begin{solution}[Basic dot product (3D)]
\[
    \vec{u} \cdot \vec{v} = (1)(4) + (-2)(1) + (3)(2) = 4 - 2 + 6 = \mathbf{8}
\]
\end{solution}

\begin{solution}[Geometric interpretation]
\begin{enumerate}[label=(\alph*)]
    \item $\theta = 0^{\circ}$, so $\cos\theta = 1$. Dot product is \textbf{positive}.
          Vectors pointing the same way are maximally aligned.
    \item $\theta = 180^{\circ}$, so $\cos\theta = -1$. Dot product is \textbf{negative}.
          The vectors point in completely opposite directions.
    \item $\theta = 90^{\circ}$, so $\cos\theta = 0$. Dot product is \textbf{zero}.
          Orthogonal vectors share no directional component.
    \item The dot product of the long and short vector is \textbf{larger}, because
          $\|\vec{a}\|\|\vec{b}\|\cos\theta$ scales with magnitude. The raw dot
          product conflates direction and magnitude --- which is why cosine similarity
          (see Exercise~1.5) is often preferred.
\end{enumerate}
\end{solution}

\begin{solution}[Dot product as similarity]
\begin{enumerate}[label=(\alph*)]
    \item $\text{``cat''} \cdot \text{``dog''} = (2)(2) + (3)(2) = 4 + 6 = \mathbf{10}$
    \item $\text{``cat''} \cdot \text{``car''} = (2)(-3) + (3)(1) = -6 + 3 = \mathbf{-3}$
    \item ``cat'' and ``dog'' are more similar by dot product (10 vs.\ $-3$). This matches
          intuition: both are animals/pets, while ``car'' is unrelated.
\end{enumerate}
\end{solution}

\begin{solution}[Normalisation and dot products]
\begin{enumerate}[label=(\alph*)]
    \item $\|\vec{p}\| = \sqrt{9+16} = \sqrt{25} = 5$. \quad $\|\vec{q}\| = \sqrt{0+25} = 5$.
    \item $\vec{p} \cdot \vec{q} = (3)(0) + (4)(5) = 20$.
    \item $\hat{\vec{p}} = \begin{bmatrix} 3/5 \\ 4/5 \end{bmatrix}$, \quad
          $\hat{\vec{q}} = \begin{bmatrix} 0 \\ 1 \end{bmatrix}$.
    \item $\hat{\vec{p}} \cdot \hat{\vec{q}} = (3/5)(0) + (4/5)(1) = 4/5 = 0.8$.

          Cosine similarity removes the effect of vector length, measuring only the
          \emph{angle} between directions. Two vectors can have a large raw dot product
          simply because they are long, even if they point in very different directions.
          Cosine similarity gives a value in $[-1, 1]$ that purely reflects directional
          alignment.
\end{enumerate}
\end{solution}

\setcounter{section}{2}
\setcounter{solution}{0}

% --- Section 2 Solutions ---
\subsection*{Section 2: Matrix Multiplication}

\begin{solution}[$2 \times 2$ matrix multiplication]
\[
    AB = \begin{bmatrix} (1)(5)+(2)(7) & (1)(6)+(2)(8) \\ (3)(5)+(4)(7) & (3)(6)+(4)(8) \end{bmatrix}
       = \begin{bmatrix} 19 & 22 \\ 43 & 50 \end{bmatrix}
\]
\end{solution}

\begin{solution}[Non-square matrix multiplication]
\begin{enumerate}[label=(\alph*)]
    \item $A$ is $2 \times 3$, $B$ is $3 \times 2$. Inner dimensions match, so $AB$
          has shape $\mathbf{2 \times 2}$.
    \item
\[
    AB = \begin{bmatrix}
        (1)(1)+(0)(0)+(2)(4) & (1)(2)+(0)(1)+(2)(0) \\
        (3)(1)+(1)(0)+(1)(4) & (3)(2)+(1)(1)+(1)(0)
    \end{bmatrix}
    = \begin{bmatrix} 9 & 2 \\ 7 & 7 \end{bmatrix}
\]
    \item Yes, $BA$ is valid: $B$ is $3 \times 2$, $A$ is $2 \times 3$, so $BA$ has
          shape $\mathbf{3 \times 3}$.
\[
    BA = \begin{bmatrix} 7 & 2 & 4 \\ 3 & 1 & 1 \\ 4 & 0 & 8 \end{bmatrix}
\]
Note: $AB \neq BA$ in general, and they don't even have the same shape here.
\end{enumerate}
\end{solution}

\begin{solution}[Shape tracking]
\begin{enumerate}[label=(\alph*)]
    \item Valid. Output shape: $\mathbf{(3 \times 2)}$.
    \item \textbf{Not valid.} Inner dimensions: 3 (from $A$) does not match 5 (from $B$).
    \item Valid. $(1 \times 6)(6 \times 1)$. Output shape: $\mathbf{(1 \times 1)}$ --- a
          scalar. This is a dot product.
    \item Valid. $(6 \times 1)(1 \times 6)$. Output shape: $\mathbf{(6 \times 6)}$ --- an
          outer product, a full matrix.
    \item $(1 \times d_{\mathrm{model}})(d_{\mathrm{model}} \times d_{\mathrm{head}})
          = \mathbf{(1 \times d_{\mathrm{head}})}$. The query vector for this token has
          $d_{\mathrm{head}}$ dimensions.
\end{enumerate}
\end{solution}

\begin{solution}[$W_V$ and $W_O$ in miniature]
\begin{enumerate}[label=(\alph*)]
    \item $\vec{v} = \vec{x} W_V
          = \begin{bmatrix} 3 & 1 \end{bmatrix}
            \begin{bmatrix} 1 & 0 \\ 0 & 2 \end{bmatrix}
          = \begin{bmatrix} 3 & 2 \end{bmatrix}$

    \item $\vec{o} = \vec{v} W_O
          = \begin{bmatrix} 3 & 2 \end{bmatrix}
            \begin{bmatrix} 0 & 1 \\ 1 & 0 \end{bmatrix}
          = \begin{bmatrix} 2 & 3 \end{bmatrix}$

    \item First, $W_V W_O =
          \begin{bmatrix} 1 & 0 \\ 0 & 2 \end{bmatrix}
          \begin{bmatrix} 0 & 1 \\ 1 & 0 \end{bmatrix}
          = \begin{bmatrix} 0 & 1 \\ 2 & 0 \end{bmatrix}$

          Then $\vec{x}(W_V W_O) =
          \begin{bmatrix} 3 & 1 \end{bmatrix}
          \begin{bmatrix} 0 & 1 \\ 2 & 0 \end{bmatrix}
          = \begin{bmatrix} 2 & 3 \end{bmatrix}$

    \item Yes --- both give $\begin{bmatrix} 2 & 3 \end{bmatrix}$. Associativity holds:
          $(\vec{x} W_V) W_O = \vec{x}(W_V W_O)$.
\end{enumerate}
\end{solution}

\begin{solution}[The OV circuit]
\begin{enumerate}[label=(\alph*)]
    \item $W_{OV} = W_V W_O
          = \begin{bmatrix} 2 & 1 \\ 0 & 3 \end{bmatrix}
            \begin{bmatrix} 1 & 0 \\ 0 & 1 \end{bmatrix}
          = \begin{bmatrix} 2 & 1 \\ 0 & 3 \end{bmatrix}$

    \item $W_O = I$ means $W_O$ does nothing. In this case $W_{OV} = W_V$ exactly.
          In practice $W_O \neq I$, and the composition $W_V W_O$ determines what
          information each head writes back to the residual stream.

    \item $\vec{x}_1 W_{OV}
          = \begin{bmatrix} 1 & 0 \end{bmatrix}
            \begin{bmatrix} 2 & 1 \\ 0 & 3 \end{bmatrix}
          = \begin{bmatrix} 2 & 1 \end{bmatrix}$

          $\vec{x}_2 W_{OV}
          = \begin{bmatrix} 0 & 1 \end{bmatrix}
            \begin{bmatrix} 2 & 1 \\ 0 & 3 \end{bmatrix}
          = \begin{bmatrix} 0 & 3 \end{bmatrix}$

          The OV matrix consistently transforms the input: the first basis direction maps
          to $[2, 1]$ and the second to $[0, 3]$. The OV circuit tells you \emph{what the
          head writes} into the residual stream when it attends to a given token ---
          independent of \emph{where} it attends.
\end{enumerate}
\end{solution}

\setcounter{section}{3}
\setcounter{solution}{0}

% --- Section 3 Solutions ---
\subsection*{Section 3: Linear Transformations}

\begin{solution}[Matrix as function]
\begin{enumerate}[label=(\alph*)]
    \item $T\begin{bmatrix}1\\0\end{bmatrix} = \begin{bmatrix}2\\0\end{bmatrix}$ ---
          stretched by 2 in the $x$-direction.
    \item $T\begin{bmatrix}0\\1\end{bmatrix} = \begin{bmatrix}0\\3\end{bmatrix}$ ---
          stretched by 3 in the $y$-direction.
    \item $T\begin{bmatrix}1\\1\end{bmatrix} = \begin{bmatrix}2\\3\end{bmatrix}$ ---
          stretched by 2 in $x$ and 3 in $y$ simultaneously.
    \item $T$ is a \textbf{scaling transformation}: it stretches the $x$-axis by factor~2
          and the $y$-axis by factor~3, independently.
\end{enumerate}
\end{solution}

\begin{solution}[Two matrices, same input]
\begin{enumerate}[label=(\alph*)]
    \item $A\vec{v} = \begin{bmatrix}1&0\\0&-1\end{bmatrix}\begin{bmatrix}3\\2\end{bmatrix}
          = \begin{bmatrix}3\\-2\end{bmatrix}$.
          $A$ \textbf{reflects} across the $x$-axis (negates the $y$-component).
    \item $B\vec{v} = \begin{bmatrix}0&1\\1&0\end{bmatrix}\begin{bmatrix}3\\2\end{bmatrix}
          = \begin{bmatrix}2\\3\end{bmatrix}$.
          $B$ \textbf{swaps} the $x$ and $y$ components --- a reflection across the
          diagonal.
    \item This is the key intuition: different weight matrices ask different questions
          about the same token vector. One head might be looking for ``is this a noun?''
          while another looks for ``is this at the start of a phrase?'' --- all operating
          on the same residual stream vector.
\end{enumerate}
\end{solution}

\begin{solution}[Matrix multiplication as function composition]
\begin{enumerate}[label=(\alph*)]
    \item $A\vec{v} = \begin{bmatrix}1&1\\0&1\end{bmatrix}\begin{bmatrix}1\\2\end{bmatrix}
          = \begin{bmatrix}3\\2\end{bmatrix}$
    \item $B(A\vec{v}) = \begin{bmatrix}2&0\\0&2\end{bmatrix}\begin{bmatrix}3\\2\end{bmatrix}
          = \begin{bmatrix}6\\4\end{bmatrix}$
    \item $C = BA = \begin{bmatrix}2&0\\0&2\end{bmatrix}\begin{bmatrix}1&1\\0&1\end{bmatrix}
          = \begin{bmatrix}2&2\\0&2\end{bmatrix}$
    \item $C\vec{v} = \begin{bmatrix}2&2\\0&2\end{bmatrix}\begin{bmatrix}1\\2\end{bmatrix}
          = \begin{bmatrix}6\\4\end{bmatrix}$
    \item Yes, both give $\begin{bmatrix}6\\4\end{bmatrix}$. Applying $A$ then $B$ is the
          same as applying the single matrix $C = BA$. This is why the OV circuit
          ($W_V W_O$) can be analysed as a single matrix.
\end{enumerate}
\end{solution}

\begin{solution}[Projections]
\begin{enumerate}[label=(\alph*)]
    \item $P\vec{v} = \begin{bmatrix}1&0\\0&0\end{bmatrix}\begin{bmatrix}4\\7\end{bmatrix}
          = \begin{bmatrix}4\\0\end{bmatrix}$. The $y$-component (7) is lost and cannot
          be recovered.
    \item $P(P\vec{v}) = P\begin{bmatrix}4\\0\end{bmatrix} = \begin{bmatrix}4\\0\end{bmatrix}$.
          Projecting again does nothing.
    \item $P^2 = PP = \begin{bmatrix}1&0\\0&0\end{bmatrix}\begin{bmatrix}1&0\\0&0\end{bmatrix}
          = \begin{bmatrix}1&0\\0&0\end{bmatrix} = P$. Geometrically: once a vector is
          on the $x$-axis, projecting it onto the $x$-axis again leaves it exactly where
          it is.
    \item The residual stream uses \textbf{addition}
          ($\vec{x}_{\mathrm{new}} = \vec{x} + \Delta\vec{x}$), not replacement.
          Nothing is discarded --- each layer adds information to the stream. A projection
          destroys information; the residual connection preserves it.
\end{enumerate}
\end{solution}

\setcounter{section}{4}
\setcounter{solution}{0}

% --- Section 4 Solutions ---
\subsection*{Section 4: Attention Scores by Hand}

\begin{solution}[Computing Q and K vectors]
Since $W_Q = I$ (the identity), query vectors equal the embeddings:
\[
    \vec{q}_1 = \begin{bmatrix}1 & 0\end{bmatrix}, \quad
    \vec{q}_2 = \begin{bmatrix}0 & 1\end{bmatrix}, \quad
    \vec{q}_3 = \begin{bmatrix}1 & 1\end{bmatrix}
\]
For keys, $W_K$ doubles the first component:
\[
    \vec{k}_1 = \begin{bmatrix}1&0\end{bmatrix}W_K = \begin{bmatrix}2&0\end{bmatrix}, \quad
    \vec{k}_2 = \begin{bmatrix}0&1\end{bmatrix}W_K = \begin{bmatrix}0&1\end{bmatrix}, \quad
    \vec{k}_3 = \begin{bmatrix}1&1\end{bmatrix}W_K = \begin{bmatrix}2&1\end{bmatrix}
\]
\end{solution}

\begin{solution}[Raw attention scores]
Computing $s_{ij} = \vec{q}_i \cdot \vec{k}_j$:
\[
    S = \begin{bmatrix} 2 & 0 & 2 \\ 0 & 1 & 1 \\ 2 & 1 & 3 \end{bmatrix}
\]
Row~3 is $[2, 1, 3]$. The highest score is $s_{33} = 3$, so token~3 most strongly
attends to \textbf{token~3} (itself).
\end{solution}

\begin{solution}[Scaling and softmax]
\textbf{Step 1:} Scale by $1/\sqrt{2} \approx 0.707$:
\[
    \tilde{S} = \begin{bmatrix} 1.41 & 0 & 1.41 \\ 0 & 0.71 & 0.71 \\ 1.41 & 0.71 & 2.12 \end{bmatrix}
\]

\begin{enumerate}[label=(\alph*)]
    \item Row~1 scaled scores: $[1.41,\ 0,\ 1.41]$.
    \item Exponentiate: $e^{1.41} \approx 4.096$, $e^0 = 1.000$, $e^{1.41} \approx 4.096$.
          Sum $= 9.192$.
\[
    A_{1,\cdot} = \left[\frac{4.096}{9.192},\ \frac{1.000}{9.192},\ \frac{4.096}{9.192}\right]
               \approx [0.446,\ 0.109,\ 0.446]
\]
Sum check: $0.446 + 0.109 + 0.446 \approx 1.0$. \checkmark

    \item Row~3 scaled scores: $[1.41,\ 0.71,\ 2.12]$.
          Exponentiate: $4.096,\ 2.034,\ 8.331$. Sum $= 14.461$.
\[
    A_{3,\cdot} \approx [0.283,\ 0.141,\ 0.576]
\]

    \item \textbf{Interpretation:}
          Token~1 splits its attention equally between itself and token~3 (both $\approx$44.6\%),
          largely ignoring token~2 ($\approx$10.9\%) --- because token~1's query has only a
          first component, and tokens~1 and~3 both have large first components in their keys.
          Token~3 attends most strongly to itself ($\approx$57.6\%), moderately to token~1
          ($\approx$28.3\%), and least to token~2 ($\approx$14.1\%).
\end{enumerate}
\end{solution}

\setcounter{section}{5}
\setcounter{solution}{0}

% --- Section 5 Solutions ---
\subsection*{Section 5: Tensor Shapes and Einsum Notation}

\begin{solution}[Basic einsum interpretation]
\begin{enumerate}[label=(\alph*)]
    \item \code{"i, i -> "}: both inputs shape \code{(4,)}, output scalar \code{()}.
          \textbf{Dot product.} Index \code{i} is contracted.
    \item \code{"i j, j k -> i k"}: \code{A} is \code{(3,4)}, \code{B} is \code{(4,5)},
          output \code{(3,5)}. \textbf{Standard matrix multiplication.}
    \item \code{"b i j, b j k -> b i k"}: \code{A} is \code{(8,3,4)}, \code{B} is
          \code{(8,4,5)}, output \code{(8,3,5)}. \textbf{Batched matrix multiplication} ---
          8 independent matrix multiplications run simultaneously.
    \item \code{"i j -> j i"}: input \code{(3,7)}, output \code{(7,3)}. \textbf{Transpose.}
\end{enumerate}
\end{solution}

\begin{solution}[Shape consistency check]
\begin{enumerate}[label=(\alph*)]
    \item \textbf{Valid --- no error.} \code{i = 3} is contracted, \code{j = 4} and
          \code{k = 5} are kept. Output: \code{(4, 5)}. Computes $A^T B$.
    \item \textbf{Error.} Index \code{j} is 4 in \code{A} but 5 in \code{B}. The inner
          dimensions don't align.
    \item \textbf{Valid --- no error.} \code{j = 3} is contracted. Output: \code{(2, 4)}.
\end{enumerate}
\end{solution}

\begin{solution}[The ARENA 1.3 einsum]
\begin{enumerate}[label=(\alph*)]
    \item \code{Q} has shape \code{(seqQ, n, h)} and \code{K} has shape
          \code{(seqK, n, h)}.
    \item Output shape: \code{(n, seqQ, seqK)}.
    \item \code{h} is contracted. Summing over \code{h} computes the \textbf{dot product}
          in the head dimension: for fixed head \code{n}, query position~$i$, and key
          position~$j$, the output is $\sum_{h} Q[i,n,h] \times K[j,n,h]$ --- exactly
          $\vec{q}_{i,n} \cdot \vec{k}_{j,n}$.
    \item Entry \code{[n, i, j]} is the \textbf{raw attention score} for head~\code{n}:
          the dot product of query at position~$i$ with key at position~$j$. After
          dividing by $\sqrt{d_{\mathrm{head}}}$ and softmaxing, this becomes the
          attention weight.
    \item \code{n} is kept because each head computes its own independent attention
          pattern. Writing \code{"seqQ n h, seqK n h -> seqQ seqK"} would sum over
          both \code{n} and \code{h}, collapsing all heads into a single score matrix
          --- destroying the multi-head structure entirely.
\end{enumerate}
\end{solution}

\begin{solution}[Writing your own einsum]
\begin{enumerate}[label=(\alph*)]
    \item \code{"m n, n p -> m p"}
    \item \code{"s d, d -> s"}
    \item \code{"m, n -> m n"} (no indices contracted --- both appear in only one input)
    \item \code{"seqK n h, n seqQ seqK -> seqQ n h"} \\
          \code{seqK} is contracted (the weighted average). \code{n}, \code{seqQ}, and
          \code{h} are kept. For each query position and head, this computes the weighted
          combination of value vectors.
\end{enumerate}
\end{solution}

\end{document}
